% !PRINTED 581
\weberpage{581}

%<<CONT
worin die $x_1, x_2, \ldots, x_n$ ganze rationale Zahlen oder ganze rationale Functionale sind, so erhält man einen Ausdruck von der Form
\weq{13}{\omega = A_0 + A_1\tau + A_2\tau^2 + \cdots + A_{n-1}\tau^{n-1},}
worin die $A_0, A_1, \ldots, A_{n-1}$ gleichfalls \gesperrt{ganze rationale Functionale} bedeuten. Nach \S.\,145, 2. wird aber in dieser Form jede ganze Zahl und jedes ganze Functional des Körpers $\Omega$ dargestellt, und wir können daher den Satz aussprechen:

3. \gesperrt{Die Potenzen}
\[
1, \tau, \tau^2, \ldots, \tau^{n-1}
\]
\gesperrt{bilden eine Basis der ganzen Functionale des Körpers $\Omega$.}

Statt der Potenzen von $\tau$ können auch die Functionen
\[
F_0, F_1, F_2, \ldots, F_{n-1}
\]
als Elemente der Basis eingeführt werden, die durch die Gleichung
\weq{14}{\frac{F(t)}{t-\tau} = t^{n-1}F_0(\tau) + t^{n-2}F_1(\tau) + \cdots + F_{n-1}(\tau)}
definirt sind, die, wie schon im \S.\,4 des ersten Bandes gezeigt ist, wenn
\[
F(t) = t^n + a_1t^{n-1} + \cdots + a_n
\]
ist, den Ausdruck haben:
\weq{15}{\begin{aligned}
F_0(t) &= 1 \\
F_1(t) &= t + a_1 \\
F_2(t) &= t^2 + a_1 t + a_2 \\
&\cdots\cdots\cdots\cdots\cdots\cdots\cdots\cdots\cdots\cdots \\
F_{n-1}(t) &= t^{n-1} + a_1 t^{n-2} + a_2 t^{n-3} + \cdots + a_{n-1},
\end{aligned}}
so dass die Potenzen $1, t, t^2, \ldots, t^{n-1}$ ohne Nenner durch $F_0, F_1, \ldots, F_{n-1}$ dargestellt werden können, und dass jede ganze Zahl oder jedes ganze Functional $\omega$ auch den Ausdruck erhält:
\weq{16}{\omega = B_0 F_0(\tau) + B_1 F_1(\tau) + \cdots + B_{n-1} F_{n-1}(\tau),}
dessen Coëfficienten $B_0, B_1, \ldots, B_{n-1}$ \gesperrt{ganze rationale Functionale} sind.

Betrachten wir jetzt die ganze Function von $t$:
\[
S\,\frac{F(t)}{(t-\tau)\,F'(\tau)},
\]
% [sic: A lowered comma is printed after "so" at the line break. It is retained although the sentence does not require this comma.]
worin $S$ das Zeichen für die in \S.\,134 erklärte Spur ist, so, ergiebt sich, dass sie den Werth 1 erhält für $t=\tau_1,\tau_2,\ldots,\tau_n$
%CONT>>
